\chapter{Annexes}
\chapminitoc

% Rappel INSA : « beaucoup de résultats et de détails sur les méthodes ->
% en annexe ». Le corps reste ~70 pages lisibles, le détail technique vit ici.

\section{Schéma électronique du conditionnement de l'encodeur (Protova1)}
% TODO : schéma RC + 74HC14, plan de câblage GP2/GP3.

\section{Plans CAO de la nouvelle base Protova1}
% TODO : mises en plan comme dans ton rapport 4A (cartouche, cotations).

\section{Scripts de lancement des véhicules}
% TODO : launch_car_wifi.sh, launch_dashboard.sh, launch_v2v_lidar.sh —
% avec 3 lignes d'explication chacun.

\section{Configurations réseau 5G / Zenoh}
% TODO : configs router/session Zenoh, commandes AT du modem.

\section{Extraits de code commentés}
% TODO : v2v_node.py (émission CAM), aeb_fusion.py (fusion), telemetry_node.py
% et un panneau React du dashboard — extraits COURTS et commentés.

\section{Mesures brutes WiFi / 5G}
% TODO : tableaux complets des campagnes de mesure.

\section{Guide d'utilisation de la plateforme}
% TODO : la « passation » aux prochains stagiaires — grosse valeur ajoutée
% (grille : « valeur des traces laissées »).
